%%%% RESUMO
%%
%% Apresentação concisa dos pontos relevantes de um texto, fornecendo uma visão rápida e clara do conteúdo e das conclusões do
%% trabalho.

\begin{resumoutfpr}%% Ambiente resumoutfpr
Reconhecer grafos que constituem uma rede de mundo pequeno é uma atividade importante
em teoria de grafos, pois essas redes possuem características de muitas redes de fenômenos
naturais e podem ser utilizadas como modelos representativos deles. Existem algumas
estratégias que podem ser adotadas para se identificar tais grafos, sendo uma delas o
coeficiente de clustering local. Apesar de ser uma métrica relativamente nova, existem vários
algoritmos já propostos com o intuito de se computar seu valor. Há uma variedade de
abordagens utilizadas para se alcançar esse objetivo, alguns através de algoritmos exatos tendo
como chave o problema de contagem de triângulos. Porém, possuindo complexidade
computacional inviável quando pensamos em grafos de redes complexas. Também há um
grande leque de opções quando pensamos em algoritmos de aproximação, utilizando das mais
variadas estratégias de amostragens. Apesar dessa miríade de algoritmos, no trabalho de Lima
et al. \cite{Lima2022} é apresentado um algoritmo que utiliza a teoria de complexidade de amostra para
não apenas aprimorar a complexidade de execução, mas também para tornar o algoritmo mais
bem adaptado para a estrutura combinatória de cada entrada. Neste presente trabalho,
pretende-se de implementar esse algoritmo e realizar uma revisão bibliográfica sobre toda a
fundamentação teórica necessária para se entender plenamente o funcionamento do algoritmo.
\end{resumoutfpr}
